Ein Elektroauto bezieht seine Energie aus einem Bleiakku mit einem Energieinhalt von \EakO. Für die elektrische Grundversorgung (z.B. Licht, Scheibenwischer, Radio, \dots) werden konstant \PgvO\ verbraucht.

\begin{abcliste}
\abc Bestimmen Sie die Reichweite des Autos bei einer konstanten Fahrgeschwindigkeit von \vfaO. Schätzen Sie dazu die benötigten Grössen sinnvoll ab.
\abc Stellen Sie mit dem Taschenrechner die Reichweite als Funktion der Geschwindigkeit graphisch dar. Bei welcher Geschwindigkeit ist die Reichweite maximal?
\abc Untersuchen Sie, wie sich die Optimierung verschiedener Parameter auf die maximale Reichweite (Geschwindigkeit und Strecke) auswirkt.
\end{abcliste}
